\section*{Capítulo \ref{ch:b3:microcanonical-ensemble} --- O ensemble microcanônico}

\begin{solution}{exo:b3:microcanonical-ensemble:1}
(a) $1, 4, 6, 4, 1$. (b) Sobre $16$: $6.25\%$, $25\%$, $37.5\%$,
$25\%$, $6.25\%$. (c) $S/k_{\text{B}} = 0$, $\ln4$, $\ln6$, $\ln4$,
$0$. (d) $n = 2$, que detém apenas $37.5\%$ da probabilidade: em
$N = 4$, o ``equilíbrio'' é mera maioria relativa --- a nitidez que
licencia a termodinâmica se compra com $N$ grande.
\end{solution}

\begin{solution}{exo:b3:microcanonical-ensemble:2}
(a) $N = 10$: $13.0$ contra $15.1$ ($14\%$ de erro); $N = 100$:
$360.5$ contra $363.7$ ($0.9\%$). (b) O termo desprezado é
$\tfrac12\ln(2\pi N)$: erro relativo $\sim\ln N/N$. (c) Aplique
Stirling aos três fatoriais. (d) Em $N = 10^{23}$, o termo descartado
vale $\sim26$ contra $N\ln2 \sim 10^{23}$: vinte e duas ordens abaixo
do termo dominante.
\end{solution}

\begin{solution}{exo:b3:microcanonical-ensemble:3}
(a) $\ln(\Omega_1\Omega_2) = \ln\Omega_1 + \ln\Omega_2$. (b)
$\Delta S = k_{\text{B}}\ln2$ --- uma partícula à qual se ofereceu o
dobro do volume. (c) $S = 156\,k_{\text{B}} = \qty{2.2e-21}{J/K}$. (d)
As entropias termodinâmicas carregam fatores de $10^{23}$: todo o
embaralhamento de todos os cassinos da Terra é invisível ao lado de um
sopro de ar morno.
\end{solution}

\begin{solution}{exo:b3:microcanonical-ensemble:4}
(a) $\Delta S = N_{\text{A}}k_{\text{B}}\ln2 = R\ln2 =
\qty{5.76}{J/K}$. (b) $2^{-N_{\text{A}}} \sim 10^{-1.8\times
10^{23}}$. (c) Mesmo amostrando configurações a cada \qty{e-10}{s}, a
espera esperada supera os \qty{4e17}{s} por um fator cujo
\emph{expoente} tem vinte e três algarismos. (d) Não impossível ---
inobservável: a lei é probabilística em princípio e absoluta em
qualquer mundo que contenha observadores de paciência finita.
\end{solution}

\begin{solution}{exo:b3:microcanonical-ensemble:5}
(a) Tome o logaritmo do produto. (b) Ao dobrar $V, N$ a $E/N$ fixo:
sem o $N!$, $S$ ganha um excesso extra do tipo $N\ln2$ --- a falha de
Gibbs. (c) Stirling sobre $\ln N!$ entrega as combinações $V/N$ e
$E/N$. (d) A célula de Planck $h^3$ fixa o zero da entropia; a
identidade das partículas fixa o $N!$ --- dois ingredientes quânticos
escondidos na termodinâmica clássica.
\end{solution}

\begin{solution}{exo:b3:microcanonical-ensemble:6}
(a) $1/T = \tfrac32 Nk_{\text{B}}/E$. (b) $P/T =
Nk_{\text{B}}/V$: a lei do gás ideal, saída da contagem. (c) A $S$
constante: $\ln(V/N) + \tfrac32\ln(E/N)$ fixo, logo $VE^{3/2}$
constante; com $E \propto T$: $TV^{2/3}$ constante, equivalente a
$PV^{5/3}$. (d) A lei dos gases, a relação energia--temperatura e o
expoente adiabático --- três pilares empíricos do volume do primeiro
ano --- são agora teoremas.
\end{solution}

\begin{solution}{exo:b3:microcanonical-ensemble:7}
(a) Simetria (blocos iguais). (b) Cada $\ln\Omega \propto
\tfrac32 N\ln E_i$: expandir em torno da divisão igual dá uma
gaussiana em $x$ de variância $\sim1/N$. (c) $x_{\text{rms}} \sim
\num{e-11}$. (d) $\Delta T/T \sim x_{\text{rms}}$: nanonanokelvins ---
as flutuações existem, e são imensuravelmente mansas na escala do
laboratório.
\end{solution}

\begin{solution}{exo:b3:microcanonical-ensemble:8}
(a) $\Omega = \binom Nn$, $S = k_{\text{B}}\ln\binom Nn$. (b)
$1/T = \partial S/\partial E = (k_{\text{B}}/\epsilon)[\ln(N-n) -
\ln n]$; resolva para $n/N$. (c) $E(T)$ sobe de $0$ até a saturação
$N\epsilon/2$; $C = \dd E/\dd T$ se anula em $T$ baixo (nenhum quantum
é acessível) e em $T$ alto (os dois níveis igualmente cheios, nada
sobra para absorver): uma corcova entre os dois --- a impressão digital
calorimétrica de qualquer população de dois níveis, usada para achar
impurezas magnéticas em sólidos. (d) É a ocupação térmica de dois
níveis --- e, com $\epsilon \to E - \mu$, ela reaparecerá literalmente
como a distribuição de Fermi--Dirac.
\end{solution}

\begin{solution}{exo:b3:microcanonical-ensemble:9}
(a) $1/T \propto \ln[(N-n)/n] < 0$ exatamente quando $n > N/2$. (b)
$T = \epsilon/[k_{\text{B}}\ln(2/3)] = -\qty{2.9}{K}$. (c) Um corpo de
$T$ negativo \emph{ganha} entropia ao perder energia, e um de $T$
positivo ganha ao receber: as duas contagens crescem quando a energia
flui do negativo $\to$ positivo --- o sistema invertido é o doador
universal, isto é, mais quente que tudo. (d) Com energia ilimitada,
$\Omega(E)$ nunca se inverte: $\partial S/\partial E$ permanece
positiva em qualquer energia finita.
\end{solution}

\begin{solution}{exo:b3:microcanonical-ensemble:10}
(a) Cada gás dobra o volume acessível: $\Delta S =
2Nk_{\text{B}}\ln2$. (b) Os $N!$ das populações fundidas absorvem
exatamente o ganho aparente: $\Delta S = 0$. (c) O aumento de entropia
é a \emph{perda de separação}: variável macroscópica alguma se moveu,
mas reseparar agora exige trabalho --- o aumento fica guardado como uma
conta futura. (d) $W_{\min} = T\Delta S =
2Nk_{\text{B}}T\ln2$; as usinas de separação isotópica (cascatas de
centrífugas de urânio, recuperação de hélio-3) pagam na prática esse
mínimo termodinâmico muitas vezes multiplicado.
\end{solution}

\begin{solution}{exo:b3:microcanonical-ensemble:11}
(a) $\lambda_T = \qty{1.6e-11}{m}$. (b) $V/N = k_{\text{B}}T/P =
\qty{4.1e-26}{m^3}$: $V/N\lambda_T^3 = \num{1.0e7}$, de modo que
$S/Nk_{\text{B}} = \ln(\num{1.0e7}) + 2.5 = 18.6$. (c) Em cheio sobre
o valor calorimétrico $18.6$. (d) A entropia absoluta contém $h$
(através de $\lambda_T$) e o $N!$: medidas de calor da era do vapor
confirmam, com três algarismos, o quantum de \omterm{def:b3:lagrangian-mechanics:action}{ação} e a identidade dos
átomos.
\end{solution}

\begin{solution}{exo:b3:microcanonical-ensemble:12}
(a) $v = (\qty{550}{nm})^3 = \qty{1.7e-19}{m^3}$: $\bar n =
\num{4.2e6}$, $\delta_{\text{rms}} = 1/\sqrt{\bar n} \approx
\num{5e-4}$. (b) Um meio perfeitamente uniforme nada espalharia para os
lados (\cref{pb:b3:perturbation-theory:1}, Parte III): essas
ondulações irredutíveis em $\sqrt N$, presentes em toda célula óptica,
são precisamente a desordem que permite ao céu existir. (c)
$\bar n = 100$: $v
\approx (\qty{16}{nm})^3$. (d) O céu é azul porque o
ar é feito de moléculas contáveis, cujos números flutuam como $\sqrt N$
--- Einstein usou exatamente isso para argumentar que as moléculas
eram reais.
\end{solution}

\begin{solution}{pb:b3:microcanonical-ensemble:1}
\textbf{1.} $n_\pm = \tfrac12(N \pm x/b)$; $\Omega =
N!/n_+!\,n_-!$.
\textbf{2.} Em $x = 0$: a cadeia livre se enovela; a extensão máxima
tem $\Omega = 1$.
\textbf{3.} Stirling sobre a binomial, expandida até segunda ordem em
$x/Nb$: a entropia gaussiana enunciada.
\textbf{4.} Com $E$ independente de $x$, qualquer puxão rumo a $x$
pequeno só pode vir da contagem: entrópico por construção.
\textbf{5.} Novelo $\sim b\sqrt N = \qty{50}{nm}$ contra $Nb =
\qty{5}{\micro m}$: a cadeia livre é cem vezes mais curta que o
próprio contorno --- amassada quase por inteiro.
\textbf{6.} O $N$, o número de elos entre ligações cruzadas: a
vulcanização encurta as cadeias efetivas, enrijecendo a rede.
\textbf{7.} $f = -T\,\partial S/\partial x =
k_{\text{B}}Tx/Nb^2$.
\textbf{8.} $k = k_{\text{B}}T/Nb^2 = \num{4.14e-21}/(\num{e4}
\times \num{2.5e-19}) = \qty{1.7e-6}{N/m}$ por cadeia.
\textbf{9.} $k \propto T$: aquecido sob carga fixa, o elástico
enrijece e \emph{se contrai}, levantando a carga; uma mola de aço
apenas amolece e cede um pouco.
\textbf{10.} $f(Nb/2) = k_{\text{B}}T/2b \approx \qty{4}{pN}$; a lei
linear falha quando $f$ se aproxima de $k_{\text{B}}T/b \approx
\qty{8}{pN}$, onde a cadeia se aproxima da extensão máxima.
\textbf{11.} $\sim\num{e14}$ cadeias em paralelo e
$\sim\num{e5}$ comprimentos de novelo em série ao longo de um
centímetro: $k_{\text{banda}} \sim k_{\text{cadeia}} \times \num{e14}/\num{e5}
\sim \num{e2}$--$\num{e3}\,\unit{N/m}$ --- o tato familiar de um
elástico.
\textbf{12.} A rigidez da borracha \emph{é} $k_{\text{B}}T$ por
configuração: mais agitação, mais recuo; a rigidez metálica é energia
de ligação, que a agitação anarmônica afrouxa.
\textbf{13.} Esticar corta a contagem configuracional; feito depressa
(sem troca de entropia com o exterior), a entropia configuracional
perdida tem de reaparecer como entropia térmica: o elástico esquenta
--- na escala de \qty{1}{K}, detectável com o lábio.
\textbf{14.} Soltar restaura as configurações; a conta térmica devolve
a diferença: ele esfria.
\textbf{15.} O peso \emph{sobe}: a contração ao aquecer é a assinatura
entrópica (os materiais elásticos por energia se dilatam).
\textbf{16.} A mola de aço se alonga um pouco (dilatação térmica) e seu
módulo cai: sinal oposto --- elasticidade de energia, não de contagem.
\textbf{17.} Os raios aquecidos se contraem, puxando a massa do aro
para fora do centro, na direção do lado frio; a gravidade exerce torque
sobre a roda desequilibrada; cada raio relaxa de novo ao girar para
longe: uma máquina térmica cuja substância de trabalho é a própria
entropia.
\textbf{18.} $k_{\text{B}}T/b = \num{4.14e-21}/\num{e-7} =
\qty{4e-14}{N} = \qty{0.04}{pN}$: piconewtons e abaixo --- exatamente a
faixa das pinças ópticas, de modo que a lei força--extensão de uma só
molécula de DNA pôde ser traçada ponto a ponto (e conferida, com os
refinamentos do item 23).
\textbf{19.} De lugar nenhum senão da contagem: na extensão $x$ há
menos \omterm{def:b3:microcanonical-ensemble:states}{microestados} que em $x - \dd x$, e a agitação térmica leva a
cadeia rumo à contagem maior; a ``força'' é $T$ vezes o gradiente de
entropia.
\textbf{20.} Em $T = 0$ nada explora configurações: a força entrópica
desaparece junto com a agitação que a alimenta --- a elasticidade seria
puramente energética, e a borracha se comportaria como um fio mole.
\textbf{21.} Gás: $P = T\,\partial S/\partial V$ com $S \ni
Nk_{\text{B}}\ln V$; cadeia: $f = -T\,\partial S/\partial x$ com a
contagem gaussiana --- as duas são a contagem empurrada contra um
\omterm{def:b3:lagrangian-mechanics:coordinates}{vínculo}, cotada a $T$.
\textbf{22.} O calor do secador entra no elástico; parte se converte em
trabalho na contração isotérmica, com as contas fechadas pela entropia
que entra e sai --- uma máquina térmica em miniatura, com a primeira
lei intacta.
\textbf{23.} A extensibilidade finita (a força real diverge perto do
esticamento máximo --- o refinamento da cadeia tipo verme) e a
cristalização induzida por deformação (o alinhamento ordena as cadeias,
dando histerese e liberação de calor além do modelo ideal).
\textbf{24.} $k_{\text{B}}T/\unit{nm} = \qty{4.1}{pN}$: a força em que
a energia térmica e os deslocamentos nanométricos se trocam em pé de
igualdade --- os motores moleculares, as proteínas que se dobram e o
DNA esticado operam todos a uma ordem de grandeza dela.
\textbf{25.} Dez mil elos cegos, contados: a \omterm{thm:b3:continuum-elasticity:hooke}{lei de Hooke} com
$k = k_{\text{B}}T/Nb^2$, calor ao esticar, pesos levantados por ar
quente e moléculas únicas obedecendo a \qty{0.04}{pN} --- a entropia,
agindo como força.
\end{solution}
